%Master File:lectures.tex

\lesson{Sequence Re-assembly}
\vspace{-2cm}
\begin{center}
  \includegraphics[height=10cm]{dna-12}
\end{center}

\keywords{Digital trees, tries, suffix trees}
%%%%%%%%%%%%%%%%%%%%%%% Next Slide %%%%%%%%%%%%%%%%%%%%%%%
\renewcommand{\Outline}{%
\begin{slide}
\section[1]{Outline}

\begin{minipage}{8cm}
  \begin{enumerate}\squeeze
    \outlineitem{Reconstruction}{rec}
    \outlineitem{Digital Trees}{dt}
    \outlineitem{Suffix Trees}{st}
    \outlineitem{Repeats}{rep}
   \end{enumerate}
\end{minipage}\hfill
\begin{minipage}{14cm}
  \vfill
  \includegraphics[width=14cm]{dna-13}
  \vfill
\end{minipage}
\end{slide}
\addtocounter{outlineitem}{1}
}

\setcounter{outlineitem}{1}
\Outline
\toptarget{firstoutline}
%%%%%%%%%%%%%%%%%%%%%%% Next Slide %%%%%%%%%%%%%%%%%%%%%%%

\begin{slide}
\section{DNA Sequencing}

\begin{PauseHighLight}
  \begin{itemize}
  \item In modern whole shotgun genome sequencing the full genome is
    broken into small pieces\pause
  \item The pieces are then read by a sequencing machine\pause
  \item This reads short sections (100--1000) bases\pause
  \item The \emph{reads} are then assembled to construct the full genome\pause
  \end{itemize}
\end{PauseHighLight}

\end{slide}

%%%%%%%%%%%%%%%%%%%%%%% Next Slide %%%%%%%%%%%%%%%%%%%%%%%

\begin{slide}
\section[-1]{Sequencing and Assembly}
\pause
\pb
\pauselevel{=1}
\begin{center}
  \multipdf[width=\linewidth]{dna}\pause
\end{center}

\end{slide}

%%%%%%%%%%%%%%%%%%%%%%% Next Slide %%%%%%%%%%%%%%%%%%%%%%%

\begin{slide}
\section{Reads and Contigs}

\begin{PauseHighLight}
  \begin{itemize}
  \item In the language of DNA reconstruction we are given a set of
    \emph{reads}\pause
  \item Our job is to reconstruct them into \emph{contigs}\pause
  \item The problem is complicated because there are errors in the reads
    (0.1-2\%)\pause
  \item Shotgun was important in speeding up the human genome project
    (there is some argument about how important)\pause
  \item What data structure should we use?\pause
  \end{itemize}
\end{PauseHighLight}

\end{slide}


%%%%%%%%%%%%%%%%%%%%%%% Next Slide %%%%%%%%%%%%%%%%%%%%%%%
\Outline % Tries
%%%%%%%%%%%%%%%%%%%%%%% Next Slide %%%%%%%%%%%%%%%%%%%%%%%

\begin{slide}
\section[-1]{Digital Trees or Tries}

\begin{PauseHighLight}
  \begin{itemize}
  \item A \emph{Trie} (pron. `try') or \emph{digital tree} is a multiway
    tree often used for storing large sets of words\pause
  \item They are trees with a possible branch for every letter of an
    alphabet\pause
  \item Their names comes from the word \textit{re\emph{trie}val}\pause
  \item Tries usually compactify the edges in the tree\pause
  \item All words end with a special letter ``\$''\pause
  \end{itemize}
\end{PauseHighLight}

\end{slide}


%%%%%%%%%%%%%%%%%%%%%%% Next Slide %%%%%%%%%%%%%%%%%%%%%%%

\begin{slide}
\section[-1]{Trie}
\pause
\pb
\pauselevel{=1}
\begin{center}
  \multipdf[width=\linewidth]{trie}\pause
\end{center}

\end{slide}


%%%%%%%%%%%%%%%%%%%%%%% Next Slide %%%%%%%%%%%%%%%%%%%%%%%

\begin{slide}
\section{General Uses of Tries}

\begin{PauseHighLight}
  \begin{itemize}
  \item Tries are yet another way of implementing sets\pause
  \item They provide quick insertion, deletion and find\pause
  \item Typically considerably quicker than binary trees and hash
    tables\pause
  \item They are particularly good for spell checkers, completion
    algorithms, longest-prefix matching, hyphenation\pause
  \item Each search finds the longest match between the words in the set
    and the query\pause
  \end{itemize}
\end{PauseHighLight}

\end{slide}

%%%%%%%%%%%%%%%%%%%%%%% Next Slide %%%%%%%%%%%%%%%%%%%%%%%

\begin{slide}
\section{Tries for Sequence Analysis}

\begin{PauseHighLight}
  \begin{itemize}
  \item Tries can be used to store DNA or proteins\pause
  \item For reconstruction of DNA tries are very useful for finding the
    next fitting section\pause
  \item They are useful when there are errors, as by searching a small
    sub-tree you can find a good match (same reason they are good for
    completion and spell checking)\pause
  \item For DNA you have a four way (5 way?) tree, which means your tree
    is deep but it doesn't waste so much memory\pause
  \end{itemize}
\end{PauseHighLight}

\end{slide}


%%%%%%%%%%%%%%%%%%%%%%% Next Slide %%%%%%%%%%%%%%%%%%%%%%%

\begin{slide}
\section[-2]{Trie for 31 Most Common English Words}
\pause
\pb
\pauselevel{=1}
\begin{center}
  \multipdf[width=0.99\linewidth]{trieTable}\pause
\end{center}
\end{slide}

%%%%%%%%%%%%%%%%%%%%%%% Next Slide %%%%%%%%%%%%%%%%%%%%%%%

\begin{slide}
\section{Disadvantage of Tries}

\begin{PauseHighLight}
  \begin{itemize}
  \item Table-based tries typically waste large amounts of memory\pause
  \item Often table-based tries are used for the first few layers, while
    lower levels use a less memory intensive data structure\pause
  \item These days memory is less of a problem so table-based tries are
    acceptable for some applications\pause
  \item There are many implementations of tries each suited to a
    particular task\pause
  \item For DNA tries are not too wasteful\pause{} (although maybe for
    proteins)\pauseb
  \end{itemize}
\end{PauseHighLight}

\end{slide}




%%%%%%%%%%%%%%%%%%%%%%% Next Slide %%%%%%%%%%%%%%%%%%%%%%%
\Outline % Suffix Trees
%%%%%%%%%%%%%%%%%%%%%%% Next Slide %%%%%%%%%%%%%%%%%%%%%%%
\begin{slide}
\section[-1]{Suffix Tree}

\begin{PauseHighLight}
  \begin{itemize}
  \item Suffix tree is a trie of all suffixes of a string\pause
  \item E.g. banana
    \begin{center}
      \includegraphics[width=\linewidth]{suffixBanana}\pause
    \end{center}
  \end{itemize}
\end{PauseHighLight}
\end{slide}

%%%%%%%%%%%%%%%%%%%%%%% Next Slide %%%%%%%%%%%%%%%%%%%%%%%

\begin{slide}
\section{Importance of Suffix Tree}

\begin{PauseHighLight}
  \begin{itemize}
  \item Fast algorithms to compute Suffix trees (and the related suffix
    array) are relative new (15-20 year old) data structures\pause
  \item They are very important for string-based algorithms\pause
  \item The classic application is in finding a match for a query
    string, $Q$, in a text, $T$\pause
  \end{itemize}
\end{PauseHighLight}

\end{slide}

%%%%%%%%%%%%%%%%%%%%%%% Next Slide %%%%%%%%%%%%%%%%%%%%%%%

\begin{slide}
\section{String Matching}

\begin{PauseHighLight}
  \begin{itemize}
  \item To find a match of a query string, $Q = an$, in a text, $T = banana$, we can
    first construct the suffix tree of the string $T$ we then simple look
    up the query, $Q$, using the trie
  \end{itemize}
    \begin{center}
      \includegraphics[width=\linewidth]{suffixBanana}\pause
    \end{center}
\end{PauseHighLight}

\end{slide}

%%%%%%%%%%%%%%%%%%%%%%% Next Slide %%%%%%%%%%%%%%%%%%%%%%%

\begin{slide}
\section{Complexity of Suffix Trees}

\begin{PauseHighLight}
  \begin{itemize}
  \item Using a regular trie for a suffix tree would typically use far
    too much memory to be useful\pause
  \item However, by using pointers to the original text it is possible
    to build a suffix tree using $O(n)$ memory where $n$ is the length
    of the text\pause
  \item Furthermore (and rather incredibly) there is a linear time
    ($O(n)$) algorithm to construct the trie\pause
  \item The algorithm is not however trivial to understand\pause
  \end{itemize}
\end{PauseHighLight}

\end{slide}

%%%%%%%%%%%%%%%%%%%%%%% Next Slide %%%%%%%%%%%%%%%%%%%%%%%

\begin{slide}
\section{Uses of Suffix Trees}

\begin{PauseHighLight}
  \begin{itemize}
  \item Suffix trees are efficient whenever it is likely that you will
    do multiple searches\pause
  \item Exact word matching is in itself a very important application\pause
  \item Suffix trees in combination with dynamic programming (which we
    look at next) can be used to do inexact matching (finding the match
    with the smallest edit distance)\pause
  \item Suffix trees get used in bioinformatics, advanced machine
    learning algorithms, \ldots\pause
  \item In case you are interested there are also prefix trees
    containing all prefixes\pause
  \end{itemize}
\end{PauseHighLight}

\end{slide}





%%%%%%%%%%%%%%%%%%%%%%% Next Slide %%%%%%%%%%%%%%%%%%%%%%%
\Outline % Repeats
%%%%%%%%%%%%%%%%%%%%%%% Next Slide %%%%%%%%%%%%%%%%%%%%%%%

\begin{slide}
\section{New Generation Sequencers}

\begin{PauseHighLight}
  \begin{itemize}
  \item The estimated cost of sequencing the human genome in 2005
    was \$$10\,000\,000$\pause
  \item To reduce the cost there was and is a drive to produce new sequencing
    machines\pause
  \item These tend to read much shorter sections of DNA (e.g.\
    20-100nt)\pause
  \item Can these be re-assembled?\pause
  \end{itemize}
\end{PauseHighLight}

\end{slide}

%%%%%%%%%%%%%%%%%%%%%%% Next Slide %%%%%%%%%%%%%%%%%%%%%%%

\begin{slide}
\section[-1]{Repeats}

\begin{PauseHighLight}
  \begin{itemize}
  \item The difficulty of assembly is caused by repeats\pause
    \begin{center}
      \includegraphics[width=\linewidth]{dna-repeat}
    \end{center}
  \item How many repeats are there in the human genome?\pause\
    (Incidentally the human genome is 3.2 billion base
    \emph{pairs})\pause
  \item You can read DNA in either directions meaning there are
    effectively 6.4 billion bases\pause
  \item This is an important question for developing new sequencing
    technologies\pause
  \end{itemize}
\end{PauseHighLight}

\end{slide}

%%%%%%%%%%%%%%%%%%%%%%% Next Slide %%%%%%%%%%%%%%%%%%%%%%%

\begin{slide}
\section[-2]{Repeats in Human Genome}

\begin{center}
  \includegraphics[width=0.7\linewidth]{nava}
\end{center}
\end{slide}

%%%%%%%%%%%%%%%%%%%%%%% Next Slide %%%%%%%%%%%%%%%%%%%%%%%

\begin{slide}
\section[-2]{Repeats Structure}

\begin{center}
  \includegraphics[width=0.7\linewidth]{nava2}
\end{center}
\end{slide}

%%%%%%%%%%%%%%%%%%%%%%% Next Slide %%%%%%%%%%%%%%%%%%%%%%%

\begin{slide}
\section{Computing Repeats}

\begin{PauseHighLight}
  \begin{itemize}
  \item A naive program would take $n^2$ operations where
    $n=6.4\times10^9$\pause
  \item If we used this we would still be waiting for the program to
    finish\pause
  \item Could not answer this question a few years ago\pause---not
    because computers weren't powerful, but because the algorithms had
    not been developed\pause
  \item However using suffix arrays counting repeats was made easy\pause
  \end{itemize}
\end{PauseHighLight}

\end{slide}

%%%%%%%%%%%%%%%%%%%%%%% Next Slide %%%%%%%%%%%%%%%%%%%%%%%

\begin{slide}
\section{Summary}

\begin{PauseHighLight}
  \begin{itemize}
  \item The incredibly swift development in algorithms has leveraged the
    advances in computer power and smart algorithms\pause
  \item Digital trees or tries, though simple can significantly speed up
    algorithms for re-assembly\pause
  \item Suffix trees are conceptually simple (digital trees of all
    suffixes)\pause
  \item They can vastly speed up operations such as matching\pause
  \item What is complicated is the algorithm to construct them
    efficiently\pause{} (though today they are widely available in programming
    languages such as python)\pause
  \end{itemize}
\end{PauseHighLight}

\end{slide}
